\section*{Bab \ref{ch:b1:ocean-circulation} --- Sirkulasi samudra}
% values: calc/ch29.py

\begin{solution}{exo:b1:ocean-circulation:1}
(a)~Pada \qty{4}{m/s} (rumus pertama): $1000\,C_D = 0.29 + 0.775 + 0.481 = 1.546$,
$T = 1.3 \times \num{1.546e-3} \times 16 = \qty{0.0322}{Pa}$. Pada \qty{7}{m/s}
(rumus kedua): $1000\,C_D = 0.60 + 0.497 = 1.097$, $T = \qty{0.0699}{Pa}$. Pada
\qty{15}{m/s}: $1000\,C_D = 1.665$, $T = \qty{0.487}{Pa}$. (b)~$0.487/0.0699 =
6.97$, dibanding $(15/7)^2 = 4.59$: koefisien seret itu sendiri bertambah dengan angin,
karena angin yang lebih kuat membangkitkan gelombang yang lebih kasar yang mencengkeram udara lebih baik.
\end{solution}

\begin{solution}{exo:b1:ocean-circulation:2}
(a)~Belahan Utara, \ang{90} ke kanan angin yang berhembus ke timur: ke selatan. (b)~Angin
berhembus ke barat daya; di sebelah kanannya adalah barat laut. (c)~Belahan Selatan, \ang{90} ke kiri angin yang berhembus ke timur:
ke utara. (d)~Angin berhembus ke utara; di sebelah kirinya adalah barat, yang berada di lepas pantai pada pantai barat. Pada kasus (d) transpor menunjuk ke lepas pantai:
air permukaan digantikan dari bawah, dan air dingin naik di sepanjang pantai.
\end{solution}

\begin{solution}{exo:b1:ocean-circulation:3}
(a)~$31.7 + 3.5 - 17.9 = \qty{17.3}{\sverdrup}$, dibanding \qty{17.1}{\sverdrup} yang diterbitkan:
rerata ketiga bagian itu menjumlah ke pembalikan rerata dalam pembulatan dan ketidakpastian ($\pm\qty{3.2}{\sverdrup}$).
(b)~Di atas \qty{1100}{m}: $16.6 + 0.4 = \qty{17.0}{\sverdrup}$ ke utara; di bawah:
$12.0 + 5.9 = \qty{17.9}{\sverdrup}$ ke selatan. Keduanya hampir seimbang, seperti yang harus terjadi
bila tidak ada air yang menumpuk di utara irisan; selisih kecil itu berada dalam ketidakpastian. (c)~$\num{17.1e6} \times \num{3.156e7} = \qty{5.40e14}{m^3}$, yaitu
\qty{5.40e5}{km^3} setahun; $\num{1.3324e9}/(\num{5.40e5}) \approx 2470$ tahun.
\end{solution}

\begin{solution}{exo:b1:ocean-circulation:4}
$D_E = 7.6 \times 10/\sqrt{\sin\varphi}$ dan $V_0 = 0.0127 \times
10/\sqrt{\sin\varphi}$. Pada \ang{15}: \qty{149}{m} dan \qty{0.250}{m/s}
(\qty{2.5}{\percent} dari angin); pada \ang{45}: \qty{90}{m} dan \qty{0.151}{m/s}
(\qty{1.5}{\percent}); pada \ang{60}: \qty{82}{m} dan \qty{0.136}{m/s}
(\qty{1.4}{\percent}). Dekat khatulistiwa $\sin\varphi \to 0$ dan kedua rumus menyimpang; lebih mendasar lagi $f \to 0$, sehingga gaya Coriolis yang dibutuhkan untuk mengimbangi gesekan lenyap, dan waktu $1/f$ yang dibutuhkan keseimbangan untuk terbentuk menjadi lebih panjang daripada angin tetap mantap. Ekman membatasi rumusnya pada $|\varphi| \geq
\ang{10}$.
\end{solution}

\begin{solution}{exo:b1:ocean-circulation:5}
(a)~$f = 2\Omega\sin\ang{11} = \qty{2.783e-5}{s^{-1}}$; $M = 0.07/(\num{2.783e-5}) =
\qty{2515}{kg.m^{-1}.s^{-1}}$, \ang{90} ke kanan angin yang berhembus ke barat:
ke utara. (b)~$Q = (2515/1027) \times \num{5e6} = \qty{1.22e7}{m^3/s}$,
\qty{12.2}{\sverdrup}. (c)~Estimasinya berada dalam ketidakpastian pengukuran langsung ($12.0 \pm 5.5$) dan dekat dengan nilai dari angin jangka panjang
($13.5 \pm 0.3$): teorema menjelaskan transpor yang teramati.
\end{solution}

\begin{solution}{exo:b1:ocean-circulation:6}
(a)~$f(\ang{2}) = \qty{5.09e-6}{s^{-1}}$; $M = 0.05/(\num{5.09e-6}) =
\qty{9824}{kg.m^{-1}.s^{-1}}$, \qty{9.57}{m^2/s}: pada \ang{2} LU ke utara (kanan angin ke barat),
pada \ang{2} LS ke selatan (kiri). (b)~Melalui sebuah strip \qty{1}{m} panjang dalam \omterm{def:b1:spherical-astronomy:geographic}{bujur}, $2 \times 9.57 = \qty{19.1}{m^3/s}$ keluar melintasi kedua \omterm{def:b1:spherical-astronomy:geographic}{lintang}, yang berjarak \qty{444.8}{km}: $w = 19.1/(\num{4.448e5}) = \qty{4.30e-5}{m/s}$,
sekitar \qty{3.7}{m} sehari. (c)~Di khatulistiwa $f = 0$ dan $|\vec M| = T/|f|$ tak terhingga: keseimbangan Ekman yang mantap tidak ada di sana, dan angka-angka di atas, yang sangat peka terhadap \omterm{def:b1:spherical-astronomy:geographic}{lintang} yang dipilih, hanya berorde besaran. Yang tetap berlaku adalah tandanya: di kedua sisinya transpor menunjuk menjauhi khatulistiwa, sehingga air permukaan menyebar dan air dingin naik di sepanjang khatulistiwa.
\end{solution}

\begin{solution}{exo:b1:ocean-circulation:7}
(a)~$f(\ang{15}) = \qty{3.775e-5}{s^{-1}}$: $M = 0.06/(\num{3.775e-5}) =
\qty{1590}{kg.m^{-1}.s^{-1}}$ ke utara; $f(\ang{40}) =
\qty{9.375e-5}{s^{-1}}$: $M = \qty{1280}{kg.m^{-1}.s^{-1}}$ ke selatan.
(b)~Antara kedua \omterm{def:b1:spherical-astronomy:geographic}{lintang}, yang berjarak \qty{2780}{km}, $\partial M_y/\partial y =
-(1590 + 1280)/(\num{2.78e6}) = \qty{-1.03e-3}{kg.m^{-2}.s^{-1}}$: dengan
\cref{prop:b1:ocean-circulation:pumping}, $w_E = \qty{-1.005e-6}{m/s}$, kecepatan ke bawah
sebesar \qty{31.7}{m} setahun. (c)~$\num{1.005e-6} \times \num{2.78e6} \times
\num{5e6} = \qty{1.40e7}{m^3/s}$, \qty{14.0}{\sverdrup} dipompa ke bawah.
\end{solution}

\begin{solution}{exo:b1:ocean-circulation:8}
(a)~$H = 1027 \times 3992 \times \num{17e6} \times 12 = \qty{8.36e14}{W}$,
\qty{0.84}{PW}. (b)~Ini adalah \qty{70}{\percent} dari \qty{1.2}{PW} yang diukur: bagian atas yang sungguhan lebih hangat
dekat permukaan daripada rerata \qty{15}{\celsius}, dan giral juga membawa panas.
(c)~$1027 \times 3992 \times \num{17e6} \times 3 =
\qty{2.1e14}{W}$, \qty{0.21}{PW}, seperempat dari (a). Hanya air yang menjadi dingin
sebelum kembali yang membawa banyak panas: transpor panas Atlantik bergantung pada
\omterm{def:b1:ocean-circulation:moc}{pembentukan air dalam}.
\end{solution}

\begin{solution}{exo:b1:ocean-circulation:9}
Misalkan $x$, $y$ dan $1 - x - y$ adalah bagian AAIW, NADW dan AABW. Keseimbangan garam dan panas, diukur dari AABW:
\[
  -0.52\,x + 0.33\,y = 34.731 - 34.67 = 0.061,\qquad
  2.9\,x + 3.6\,y = 2.78 - 0.4 = 2.38 .
\]
Determinannya adalah $-0.52 \times 3.6 - 0.33 \times 2.9 = -2.829$, dan aturan Cramer memberi $x = (0.061 \times 3.6 - 0.33 \times 2.38)/(-2.829) = 0.20$
dan $y = (-0.52
\times 2.38 - 2.9 \times 0.061)/(-2.829) = 0.50$: sampelnya \qty{20}{\percent} AAIW, \qty{50}{\percent} NADW dan \qty{30}{\percent} AABW,
titik di dalam segitiga ketiga air.
\end{solution}

\begin{solution}{exo:b1:ocean-circulation:10}
(a)~Dengan $W = u + iv = V_0e^{i\pi/4}e^{(1+i)az}$, yang bagian riil dan imajinernya adalah rumus proposisi, $W'' = (1+i)^2a^2W = 2ia^2W$. Kedua persamaan berpadu menjadi
$A_z(u'' + iv'') = -fv + ifu = if(u + iv)$, yaitu $A_zW'' = ifW$, yang berlaku karena $2ia^2A_z = if$ untuk $a^2 = f/(2A_z)$. (b)~$W'(0) =
(1+i)aV_0e^{i\pi/4} = \sqrt2\,aV_0\,e^{i\pi/2} = i\sqrt2\,aV_0$: tegangan permukaan $\rho A_zW'(0)$ menunjuk ke utara, dengan besar $\rho A_z\sqrt2\,aV_0 = \rho A_z\sqrt2
\sqrt{f/(2A_z)}\;T/(\rho\sqrt{fA_z}) = T$. (c)~$\int_{-\infty}^0 W\,\mathrm dz =
V_0e^{i\pi/4}/((1+i)a) = V_0/(\sqrt2\,a)$, bilangan riil: transpor tidak memiliki bagian ke utara, dan bagian timurnya adalah $\rho V_0/(\sqrt2\,a) =
T/\sqrt{fA_z}\cdot\sqrt{A_z/f} = T/f$. Ini adalah teorema: besar $T/f$, \ang{90} ke kanan angin yang berhembus ke utara.
\end{solution}

\begin{solution}{exo:b1:ocean-circulation:11}
(a)~Langkah ke utara $\mathrm dy = a_\oplus\,\mathrm d\varphi$ mengubah $f = 2\Omega
\sin\varphi$ sebesar $2\Omega\cos\varphi\,\mathrm d\varphi$, sehingga $\beta = 2\Omega
\cos\varphi/a_\oplus = 2 \times \num{7.292e-5} \times \cos\ang{30}/(\num{6.371e6}) =
\qty{1.98e-11}{m^{-1}.s^{-1}}$. (b)~$R/\beta = \num{e-6}/(\num{1.98e-11}) =
\qty{5.04e4}{m}$, sekitar \qty{50}{km}. (c)~Interior: $\num{30e6}/(\num{5e6} \times
1000) = \qty{6.0e-3}{m/s}$, \qty{6}{mm/s}; arus batas: $\num{30e6}/(\num{5.04e4}
\times 1000) = \qty{0.60}{m/s}$, seratus kali lebih cepat, rasio lebarnya. Arus Teluk, pada \qty{1.8}{m/s}, lebih cepat lagi: arus sungguhan lebih sempit daripada perkiraan ini, alirannya terkonsentrasi dekat permukaan, dan resirkulasinya menambah transpornya.
\end{solution}

\begin{solution}{exo:b1:ocean-circulation:12}
(a)~Di mana pun laut lebih dalam dari \qty{1000}{m}, kolom di atas \qty{1000}{m} menampung paling banyak \qty{1}{km} air, dan di tempat lain lebih sedikit: volume di atas \qty{1000}{m} paling banyak $\num{361.84e6} \times 1 = \qty{3.618e8}{km^3}$, sehingga volume di bawahnya setidaknya $\num{1.3324e9} - \num{3.618e8} = \qty{9.71e8}{km^3}$.
(b)~$\num{26e6} \times \num{3.156e7} = \qty{8.21e14}{m^3}$ setahun,
\qty{8.21e5}{km^3}; $\num{9.71e8}/(\num{8.21e5}) \approx 1200$ tahun, orde yang sama
dengan seribu tahun konveyor. (c)~$(1/2)^{25/12.3} = 0.244$, sekitar seperempat;
$\num{e-3} \times 25 \times \num{3.156e7} = \qty{7.9e5}{m}$, sekitar \qty{790}{km}.
Setelah 25 tahun tritium telah menyebar tidak lebih jauh daripada yang dibawa oleh hanyut lambat,
sementara seperempatnya masih tersisa untuk dideteksi.
\end{solution}

\begin{solution}{pb:b1:ocean-circulation:1}
\begin{enumerate}
  \item Angin dari utara berhembus ke selatan, sepanjang pantai.
  \item $1000\,C_D = 0.60 + 0.071 \times 8 = 1.168$: $C_D = \num{1.168e-3}$.
  \item $T = 1.3 \times \num{1.168e-3} \times 64 = \qty{0.0972}{Pa}$.
  \item $C_D = \num{1.452e-3}$ dan $T = \qty{0.272}{Pa}$: tegangan tumbuh sebesar $2.80$,
        lebih dari $(12/8)^2 = 2.25$, karena $C_D$ itu sendiri bertambah dengan angin.
  \item $\Omega = \qty{7.292e-5}{rad.s^{-1}}$; $f = 2\Omega\sin\ang{38} =
        \qty{8.979e-5}{s^{-1}}$.
  \item $2\pi/f = \qty{7.00e4}{s}$, \qty{19.4}{h}. Gaya Coriolis memutar arus melalui
        satu lingkaran penuh dalam waktu ini; sebelum angin berhembus sekitar selama itu,
        air berosilasi alih-alih menetap ke dalam keseimbangan mantap teorema.
  \item \ang{90} ke kanan angin yang berhembus ke selatan: ke barat, laut terbuka.
  \item $M = T/f = 0.0972/(\num{8.979e-5}) = \qty{1082}{kg.m^{-1}.s^{-1}}$.
  \item $1082/1027 = \qty{1.054}{m^2/s}$ untuk setiap meter pantai.
  \item $D_E = 7.6 \times 8/\sqrt{\sin\ang{38}} = \qty{77.5}{m}$.
  \item $1.054/77.5 = \qty{0.0136}{m/s}$, \qty{1.36}{cm/s}.
  \item $V_0 = 0.0127 \times 8/\sqrt{\sin\ang{38}} = \qty{0.129}{m/s}$, \ang{45} ke
        kanan angin: ke barat daya. Pelampung memberi sekitar setengah nilai Ekman,
        sekitar \qty{0.07}{m/s}.
  \item Air yang meninggalkan pantai dalam \omterm{def:b1:ocean-circulation:ekman}{lapisan Ekman} tidak dapat digantikan melintasi pantai, maupun dari atas: ia digantikan oleh air yang naik dari bawah lapisan
        dekat pantai, seperti dalam irisan melintang yang digambar dalam
        \cref{sec:b1:ocean-circulation:upwelling}; isoterm membelok naik ke arah pantai.
  \item $1.054 \times 1000 \times 86400 = \qty{9.11e7}{m^3}$ sehari per kilometer
        pantai.
  \item $w = 1.054/(\num{2e4}) = \qty{5.27e-5}{m/s}$, \qty{4.55}{m} sehari.
  \item $77.5/4.55 = 17$ hari.
  \item Udara didinginkan dari bawah oleh air dingin: lapisan dekat permukaan menjadi
        dingin dan stabil, kelembapannya dapat mengembun menjadi kabut dan stratus rendah,
        dan tanpa konveksi terdapat sedikit hujan di pantai.
  \item Air yang di-upwell kaya nutrien; nutrien itu memberi makan fitoplankton, yang memberi makan zooplankton dan ikan.
  \item Angin yang berhembus ke utara menggerakkan transpor \ang{90} ke kanannya, ke timur: air permukaan menumpuk di pantai dan tenggelam (\omterm{def:b1:ocean-circulation:upwelling}{downwelling}); \omterm{def:b1:ocean-circulation:upwelling}{upwelling} berhenti dan air permukaan yang lebih hangat kembali ke pantai.
  \item $f(\ang{20}) = \qty{4.988e-5}{s^{-1}}$: $M = \qty{1948}{kg.m^{-1}.s^{-1}}$,
        $\sin\ang{38}/\sin\ang{20} = 1.80$ kali lebih banyak.
  \item Di Belahan Selatan transpor berada \ang{90} ke kiri angin. Untuk mendorongnya ke barat, lepas pantai, angin harus berhembus ke utara, dari selatan: dalam kedua kasus angin berhembus sepanjang pantai ke arah khatulistiwa.
  \item $1000\,C_D = 0.29 + 0.62 + 0.308 = 1.218$; $T = \qty{0.0396}{Pa}$; $M =
        \qty{441}{kg.m^{-1}.s^{-1}}$; rasionya $0.407$ untuk keduanya.
  \item Per meter pantai, $M c_p\,\Delta T = 1082 \times 3992 \times 6 =
        \qty{2.59e7}{W}$, \qty{25.9}{MW}; sepanjang \qty{1000}{km}, \qty{2.59e13}{W},
        \qty{25.9}{TW}.
  \item $\num{31.7e6}/1.054 = \qty{3.01e7}{m}$, sekitar \qty{30000}{km}: \omterm{def:b1:ocean-circulation:upwelling}{upwelling pantai} memindahkan air jauh lebih sedikit daripada \omterm{def:b1:ocean-circulation:wbc}{arus batas barat}.
  \item $Q = 1.054 \times \num{e6} = \qty{1.05e6}{m^3/s}$: transpor \omterm{def:b1:ocean-circulation:upwelling}{upwelling} sepanjang \qty{1000}{km} pantai adalah \qty{1.05}{\sverdrup}.
\end{enumerate}
\end{solution}
